\chapter{Sustainability Analysis and Ethical Implications}

This chapter follows the ETSETB sustainability guide structure: first a
sustainability matrix, then the ethical implications, and finally the relation
with the Sustainable Development Goals. The analysis separates the development
of the BT from the hypothetical execution of the project, because this thesis
does not manufacture a physical product. Its deployable result is a software
pipeline for faster CKM prediction.

\section{Sustainability matrix}

\begin{table}[ht]
\centering
\caption{Sustainability matrix summary for the thesis.}
\label{tab:sustainability_matrix}
\begin{tabularx}{\textwidth}{lXXX}
\toprule
Perspective & Development of BT & Project execution & Risks and limitations \\
\midrule
Environmental &
GPU training, laptop use, storage, and writing. &
Reduced need for repeated ray-tracing simulations. &
Extra retraining, larger models, or poor reuse of cached priors could increase footprint. \\
Economic &
Student and supervision time, plus shared HPC use. &
Lower planning cost if the surrogate replaces repeated simulations. &
Real deployment needs field validation and maintenance. \\
Social &
Skills gained in wireless modelling, ML, and scientific reporting. &
Faster UAV coverage planning for temporary or emergency networks. &
Bias toward represented city morphologies and risk of overclaiming simulated results. \\
\bottomrule
\end{tabularx}
\end{table}

\subsection{Environmental impact}

\textbf{Development of the BT.} The main environmental cost of the academic
work was electricity for computation. The final Try~80 line used four RTX~2080
Ti-class GPUs for approximately \SI{12}{h}, consuming about \SI{18}{kWh}.
The final Try~76--79 stage required roughly \SI{18}{h} of similar GPU time,
adding about \SI{27}{kWh}. The complete exploratory project exceeded
\SI{400}{GPU-h}; at the same approximation of \SI{1.5}{kW} per four-GPU node,
this corresponds to roughly \SI{150}{kWh} and approximately
\SI{30}{kgCO2e} at an illustrative factor of 0.2~kgCO2e/kWh. This
exploratory footprint is a one-time research overhead, not the marginal cost
of using the final system in deployment.

Several decisions reduced avoidable impact: using city-holdout validation
instead of repeated full retraining for every diagnostic, reusing calibration
artifacts from Try 78 and Try 79 in Try 80, plotting from saved histories, and
keeping final priors frozen rather than recomputing them blindly in every
experiment.

\textbf{Project execution.} In operation, a trained CKM surrogate can reduce
the number of ray-tracing simulations required during planning. A ray tracer is
still needed to generate reliable training data and to validate new regimes, 
but this is a one time thing. And now this model can generalize without needing
new simulations for 6G enabled 7.125GHz UAV enabled networks, which is the main future use case.
The surrogate can also be used for fast what-if analyses during planning: 
it is not a replacement for ray tracing in every scenario, but it can evaluate
many UAV heights or urban alternatives at much lower marginal energy cost.


\textbf{Risks and limitations.} The environmental benefit depends on reuse. If
the model is retrained from scratch for every new city or frequency, the energy
advantage becomes weaker, but this shouldn't be the case unless very different cities are considered. 
The analysis also does not include the embodied carbon of the GPUs or the data-centre infrastructure,
because the project used existing shared equipment.

\subsection{Economic impact}

\textbf{Development of the BT.} Table~\ref{tab:sustainability_budget} gives an
indicative professional-cost view. It is not a billed project budget, but it
answers the guide's request to quantify human and material resources.

\begin{table}[htbp]
\centering
\footnotesize
\renewcommand{\arraystretch}{1.1}
\caption{Indicative economic cost of the BT development.}
\label{tab:sustainability_budget}
\begin{tabularx}{\textwidth}{>{\raggedright\arraybackslash}Xrrr>{\raggedright\arraybackslash}p{3cm}}
\toprule
Task or resource & Quantity & Unit cost & Cost & Notes \\
\midrule
Literature review and proposal     & 60\,h         & 12\,EUR/h        &    720\,EUR & Initial scope and SOA \\
Dataset and baseline implementation & 90\,h        & 12\,EUR/h        & 1\,080\,EUR & HDF5, U-Net/cGAN, masking \\
Experiment design and debugging    & 180\,h        & 12\,EUR/h        & 2\,160\,EUR & Tries, priors, experts, evaluation \\
Writing and figures                & 120\,h        & 12\,EUR/h        & 1\,440\,EUR & Thesis text and diagrams \\
Supervision meetings and his time for dataset generation from simulation           & 50\,h         & 40\,EUR/h        & 2\,000\,EUR & Advisor time \\
GPU compute, final Try\,80         & 48\,GPU-h     & 0.25\,EUR/GPU-h  &    12\,EUR  & 4$\times$ RTX\,2080\,Ti, 12\,h \\
GPU compute, Tries 76--79          & 72\,GPU-h     & 0.25\,EUR/GPU-h  &    18\,EUR  & 4$\times$ RTX\,2080\,Ti, 18\,h \\
Exploratory compute                & $\approx$400\,GPU-h & 0.25\,EUR/GPU-h & $\approx$100\,EUR & Research iterations \\
Electricity                        & ---           & incl.\ in GPU-h  &     0\,EUR  & Covered by GPU rental rate \\
Storage and code hosting           & ---           & free             &     0\,EUR  & GitHub \\
\midrule
\textbf{Total indicative cost}     &               &                  & \textbf{$\approx$7\,530\,EUR} & \\
\bottomrule
\end{tabularx}
\end{table}

\textbf{Project execution.} The project is economically attractive if the
trained surrogate reduces repeated ray-tracing runs during network planning.
The main future costs would be dataset generation for new frequencies or
regions, occasional recalibration with field data, model maintenance, and
quality assurance before operational use.

\textbf{Risks and limitations.} The largest economic risk is a simulator-to-real
gap. If real measurements show that the ray-traced dataset does not transfer,
additional field campaigns and recalibration would be necessary. Another risk
is scope growth: making a frequency-general outdoor CKM model requires new
simulations, larger storage, and more validation.

\subsection{Social impact}

\textbf{Development of the BT.} The work had a positive educational impact: it
required combining propagation theory, machine learning, HPC workflows,
scientific writing, and reproducible evaluation. The thesis also required
continuous critical review of failed approaches, which is a useful professional
practice. Inclusive language is used where applicable, and the work does not
assign technical roles or expected users by gender.

\textbf{Project execution.} Faster CKM prediction can help plan temporary UAV
coverage for public events, emergency response, rural extensions, or
infrastructure restoration. The direct beneficiaries would be network planners
and, indirectly, users who receive more reliable wireless coverage.

\textbf{Risks and limitations.} The project could negatively affect users if
coverage decisions were made from overconfident simulated predictions without
field validation. There is also a representativeness risk: the 66-city dataset
is broad, but no dataset covers every urban morphology, construction material,
or socioeconomic context. Future systems should monitor errors by city type,
height bin, and LoS/NLoS regime rather than using one global RMSE as the only
quality indicator.

\section{Ethical implications}

The work addresses the need for faster and more transparent radio-map
prediction in UAV-enabled networks. That need was initially defined by the
project proposal and refined through experiment results: pure deep learning was
not enough, so the final method keeps interpretable priors and reports
regime-level errors.

The main ethical obligation is honesty about validity. The dataset is simulated
through ray tracing, so the thesis must not claim real-world deployment
performance without field measurements. The final Try 80 results are also left
as placeholders until training finishes and a frozen checkpoint is selected.

The project aligns with professional engineering ethics by documenting
assumptions, avoiding data leakage between cities, reporting limitations, and
making the evaluation contract explicit. A future operational system should
also include human review before coverage decisions affect public-service or
emergency deployments.

\section{Relationship with the Sustainable Development Goals}

The thesis is most closely aligned with:

\begin{itemize}
  \item \textbf{SDG 9: Industry, Innovation and Infrastructure.} The project
  develops a method for faster planning and evaluation of wireless
  infrastructure.
  \item \textbf{SDG 11: Sustainable Cities and Communities.} Better radio-map
  prediction can support temporary and resilient urban connectivity. With less needed resources to plan wireless infrastructure.
  \item \textbf{SDG 13: Climate Action.} The method can reduce repeated
  high electricity (hence carbon) costs of simulation once a validated surrogate is available, although this
  benefit depends on responsible reuse and avoiding unnecessary retraining.
\end{itemize}

The contribution is indirect: this thesis does not deploy infrastructure by
itself, but it can reduce planning cost and help future UAV network tools make
coverage decisions faster and with clearer uncertainty reporting.
